\documentclass[11pt]{article}
\usepackage[margin=1in]{geometry}
\usepackage{amsmath, amssymb}
\usepackage{graphicx}
\usepackage{booktabs}
\usepackage{caption}
\usepackage[hidelinks]{hyperref}

\graphicspath{{../outputs/validation/}}

\title{Multi-taper spectral estimation for PUV processing:\\
       method comparison and recommendation}
\author{Holden Leslie-Bole}
\date{April 2026}

\begin{document}
\maketitle

\section{Background}

\paragraph{Note (May 2026):} This writeup documents the spectral-
estimator comparison conducted in April 2026 on the legacy 17-minute
segmentation. The conclusion was that full multi-taper ($NW=4$,
$K=7$) is the best estimator for PUV spectra. In May 2026 the L2
default segment length was further changed from 17 minutes to 1
hour to match MOP cadence and to give the longer averaging window
that segment-mean velocity diagnostics benefit from (see
mean-flow validation methods, Section~\ref{sec:mean_flow_validation}
in the methods writeup, and within-hour stationarity audit). The
production pipeline therefore now combines a 1-hour segment with
full MTM (NW=4, K=7), giving $\Delta f \approx \SI{2.78e-4}{\hertz}$
and $\sim$14 effective DOF per frequency bin per segment. The
analyses below remain valid as a comparison of estimators at fixed
segment length.

The PUV L2 spectral processor computes surface elevation spectra,
directional coefficients, and derived wave parameters from 17-minute
segments (2048 samples at 2~Hz). The original implementation uses
Welch's method with a Hanning window, \texttt{nfft}~=~256, and 50\%
overlap, giving a frequency resolution of $\Delta f = 0.0078$~Hz.
This places approximately 3 frequency bins across a typical Southern
California swell peak at 0.06--0.08~Hz.

A multi-taper approach (Thomson, 1982) may be beneficial for
two reasons: (1) DPSS tapers are optimal for minimizing spectral
leakage, and (2) the variance reduction from multiple tapers enables
longer FFT segments without sacrificing estimate stability, improving
frequency resolution.

\section{Methods}

Three spectral estimation methods were tested on TBR23 MOP580\_7m
(8.3~m depth, 4795 valid segments):

\begin{enumerate}
  \item \textbf{Hanning/Welch} (baseline): Hanning window,
    \texttt{nfft}~=~256, 50\% overlap. Approximately 15 sub-segments
    averaged per 17-minute record. $\Delta f = 0.0078$~Hz.

  \item \textbf{Multi-taper hybrid}: DPSS tapers ($NW = 4$, $K = 7$
    tapers) applied within Welch's framework using the same
    \texttt{nfft}~=~256 sub-segments with 50\% overlap. Averages over
    both tapers and sub-segments ($\sim$105 estimates per frequency
    bin). $\Delta f = 0.0078$~Hz.

  \item \textbf{Full multi-taper}: DPSS tapers ($NW = 4$, $K = 7$)
    applied to the full 2048-sample segment. No sub-segment averaging;
    variance controlled entirely by the 7 tapers.
    $\Delta f = 0.00098$~Hz ($\sim$24 bins across the swell peak).
\end{enumerate}

All three methods use the same tapered-FFT machinery for both auto-
and cross-spectra, preserving consistency between $S_{pp}$, $S_{uu}$,
$S_{vv}$, $S_{pu}$, $S_{pv}$, and $S_{uv}$. Normalization was
verified against Parseval's theorem using white noise (integral of PSD
$= \sigma^2$ to within 1\% for all three methods). Downstream
processing (pressure correction, bulk parameters, directional
coefficients, radiation stress) is identical.

\section{Spectral peakedness metric}

The Goda peakedness parameter $Q_p$ (Goda, 1970) quantifies how
concentrated spectral energy is around the peak frequency:
\begin{equation}
  Q_p = \frac{2}{m_0^2} \int_0^\infty f \, S_\eta^2(f) \, df
\end{equation}
where $m_0 = \int S_\eta(f)\,df$ is the zeroth spectral moment (total
variance). $Q_p = 1$ for a flat (white noise) spectrum; narrow-banded
swell produces $Q_p \sim 2$--4. A higher $Q_p$ indicates a sharper,
more peaked spectrum. Because $Q_p$ depends on $S_\eta^2$, it is
sensitive to the spectral estimator: a smoother PSD (e.g., Welch with
coarse $\Delta f$) will produce lower $Q_p$ than a well-resolved
estimate of the same wave field.

\section{Results}

\subsection{Bulk wave parameters}

Table~\ref{tab:summary} compares median bulk parameters across all
4795 valid segments. $H_s$ is essentially identical: the full
multi-taper differs from Welch by 0.05~cm (0.08\%). Energy flux
$E_f$ and radiation stress $S_{xx}$ are likewise unchanged. This
confirms that integral quantities are insensitive to the spectral
estimator, as expected from Parseval's theorem.

\begin{table}[h]
  \centering
  \caption{Median bulk parameters for TBR23 MOP580\_7m (4795 segments).}
  \label{tab:summary}
  \begin{tabular}{@{}lccc@{}}
    \toprule
    Statistic & Welch & MTM hybrid & MTM full \\
    \midrule
    $H_s$ (m)             & 0.615 & 0.619 & 0.615 \\
    $T_p$ (s)             & 12.80 & 10.67 & 13.30 \\
    $E_f$ (W/m)           & 1560  & 1568  & 1561  \\
    $S_{xx}$ (N/m)        & 240   & 231   & 241   \\
    Z-test (SS)           & 0.789 & 0.806 & 0.791 \\
    Spread (SS, deg)      & 17.6  & 19.1  & 16.6  \\
    \bottomrule
  \end{tabular}
\end{table}

Peak period $T_p$ is the one parameter that differs substantially.
The Welch estimate (12.80~s) is quantized to $\Delta f = 0.0078$~Hz
steps, meaning $T_p$ for 14-second swell can only take values of
12.8, 14.2, or 15.9~s---jumps of more than a second. The full
multi-taper ($\Delta f = 0.00098$~Hz) resolves $T_p$ to 13.30~s. The
hybrid, despite using the same \texttt{nfft}~=~256, shifts $T_p$ to
10.67~s because the DPSS taper shape alters the relative heights of
adjacent spectral bins at coarse resolution.

$H_s$ bias relative to Welch is +0.4~cm for the hybrid and +0.05~cm
for the full multi-taper (Fig.~\ref{fig:bulk}).

\subsection{Spectral shape}

Figure~\ref{fig:spectra} shows the time-averaged surface elevation
spectrum. The Welch estimate smears the swell peak across 3 bins. The
full multi-taper resolves a sharp, narrow peak with fine structure
that Welch cannot distinguish. The hybrid is intermediate---same
frequency resolution as Welch but slightly different taper response.

The single-segment comparison (Fig.~\ref{fig:single}) is more
striking: at the peak segment ($H_s = 1.50$~m), the full multi-taper
reveals a swell peak 2--3$\times$ taller and half the width of the
Welch estimate. This is not additional energy---it is the same energy
concentrated into a resolved peak rather than spread across adjacent
bins.

\subsection{Directional estimates}

Swell-band directional spread narrows from 17.6$^\circ$ (Welch) to
16.6$^\circ$ (full MTM). This $\sim$1$^\circ$ decrease is consistent
with reduced cross-frequency averaging: finer resolution means each
frequency bin samples a narrower directional window, reducing
apparent spread. The hybrid gives wider spread (19.1$^\circ$), likely
because DPSS tapers applied to short 256-sample sub-segments
introduce more cross-frequency leakage than Hanning at this segment
length.

\subsection{Z-test}

The pressure-velocity consistency ratio is unchanged across methods
(median $\sim$0.79 in the SS band for all three;
Fig.~\ref{fig:ztest}). The systematic departure from 1.0 is a
property of the data (sensor geometry, non-hydrostatic effects), not
the spectral estimator.

\section{Implications for MOP comparison}

The four-hypothesis analysis of the PUV--MOP spectral excess
attributed the positive $H_s$ bias to MOP spectral peak broadening
(H4). The MOP validation with full multi-taper processing is below
to test whether the hypothesis results change. Table~\ref{tab:hyp}
shows the comparison.

\begin{table}[h]
  \centering
  \caption{MOP validation results: Welch vs.\ full multi-taper
    (TBR23 MOP580\_7m).}
  \label{tab:hyp}
  \begin{tabular}{@{}lcc@{}}
    \toprule
    Metric & Welch & MTM full \\
    \midrule
    Median $Q_p$ (PUV)            & 1.88  & 2.15  \\
    Median $Q_p$ (MOP)            & 1.75  & 1.74  \\
    $Q_p$ ratio (PUV/MOP)         & 1.07  & 1.24  \\
    H4: $R(\Delta Q_p,\;\text{Hs ratio})$ & 0.443 & 0.463 \\
    H4: $R^2$                     & 0.197 & 0.214 \\
    \midrule
    H3: Low-swell ratio (PUV/MOP) & 1.048 & 1.060 \\
    PUV spread (deg)              & 15.2  & 14.3  \\
    MOP spread (deg)              & 9.8   & 9.8   \\
    \bottomrule
  \end{tabular}
\end{table}

The Welch/Hanning window estimator was broadening the PUV swell peak, partially
masking the true peakedness difference. With full multi-taper, PUV
$Q_p$ increases from 1.88 to 2.15 (14\% sharper) while MOP $Q_p$
is unchanged (1.74--1.75). The $Q_p$ ratio nearly doubles from 1.07
to 1.24: the PUV spectrum is 24\% more peaked than MOP, compared to
the 7\% difference that Welch estimated. The H4 correlation
strengthens from $R = 0.443$ to $R = 0.463$ ($R^2$: 0.197 to
0.214).

The other hypotheses are unaffected. H3 (bound long waves) remains
ruled out: the low-swell ratio shifts marginally from 1.048 to
1.060. H2 (directional narrowing) remains ruled out: PUV spread
narrows slightly from 15.2$^\circ$ to 14.3$^\circ$ with MTM, but
is still wider than MOP (9.8$^\circ$).

The multi-taper result does not change which hypothesis is supported.
H4 was and remains the only viable explanation. What it does
change is the magnitude of the effect. The MOP model broadens the
spectral peak by $\sim$24\% relative to the true spectrum, not
$\sim$7\% as originally estimated with Welch processing.

\section{Field validation against an independent prior pipeline}
\label{sec:ruby2d}

The internal MOP-comparison results in the previous section establish
that switching to multi-taper changes spectral shape diagnostics in
the expected direction without changing bulk parameters. To verify
that the new estimator does not introduce biases relative to existing
processing conventions in the lab, we ran a head-to-head test against
an independent prior pipeline
(\texttt{beachpeeps/PUV\_Processing}) on the Ruby2D deployment.

\subsection{Setup}

We reprocessed the Ruby2D MOP582\_6m record (Torrey Pines, October 2021
to February 2022, $\sim$3,100~hours) by adapting the legacy archived
L1 file directly into the new L2 processor. The L1 struct from the prior pipeline
shares its key fields (\texttt{time}, \texttt{P}, \texttt{BuoyCoord},
\texttt{fs}) with the new format, so the only adapter required was setting
\texttt{label}, \texttt{deploymentName}, and \texttt{doffp}.

The two pipelines differ in three relevant respects:
\begin{enumerate}
  \item \emph{Spectral estimator}: the new pipeline uses full multi-taper
    ($NW = 4$, $K = 7$ DPSS tapers, $\Delta f = 0.000976$~Hz);
    the legacy estimator is a single Hanning periodogram on the full 60-minute segment
    ($\Delta f = 0.000278$~Hz).
  \item \emph{Window consistency}: the new pipeline uses the same DPSS tapers for
    both auto- and cross-spectra; the legacy code uses Hanning for auto-spectra
    and Hamming for cross-spectra.
  \item \emph{Pressure correction}: the new pipeline zeros frequencies where
    $K_p < 0.1$; the legacy code caps $K_p^2 \le 10$ with no SNR-based cutoff.
\end{enumerate}

The legacy 60-minute segments were matched to the new 17-minute segments by
midpoint nearest neighbor with a 30-minute tolerance, yielding 2,322
matched pairs out of the 3,095 total legacy segments (75\%; the remaining
25\% are hours that the new pipeline rejects under
\texttt{nanMaxFrac} = 0.10 but that the legacy pipeline retains by
zero-padding NaN gaps).

\subsection{Bulk parameter agreement}

\begin{table}[h]
  \centering
  \caption{Ruby2D MOP582\_6m: bulk wave parameter agreement, the new
    multi-taper pipeline vs.\ the legacy archived L2. $T_p$ statistics
    are computed after restricting both estimates to the SS band
    (4--25~s); without this filter, three segments on January 15,
    2022 in the prior pipeline report $T_p$ at the IG/DC bin
    (1200, 1800, 3600~s) because the prior pipeline's peak picker is
    not band-limited.}
  \label{tab:ruby2d}
  \begin{tabular}{@{}lccccc@{}}
    \toprule
    Metric & Legacy (med.) & Multi-taper (med.) & Bias (M$-$L) & RMS & $R^2$ \\
    \midrule
    $H_s$ (m)             & 0.771  & 0.757  & $-0.009$ & 0.053 & 0.981 \\
    $T_p$ (s)             & 13.48  & 13.30  & $-0.02$  & 1.48  & 0.761 \\
    Direction SS ($^\circ$)    & $-7.3$ & $-7.2$ & $+0.2$   & 1.2   & 0.927 \\
    Spread SS ($^\circ$)       & 14.7   & 15.8   & $+1.0$   & 1.7   & 0.877 \\
    \bottomrule
  \end{tabular}
\end{table}

The agreement is excellent (Table~\ref{tab:ruby2d}). $H_s$ matches
to 5~cm RMS with a $-9$~mm median bias; direction matches to
$1.2^\circ$~RMS. The spread is $1^\circ$ wider in the new pipeline on
average, the largest systematic offset observed and well within
segment-to-segment scatter. The window-mismatch issue in the prior
pipeline (auto-Hanning, cross-Hamming) does not produce the
order-20\% directional errors that we initially feared --- at Ruby2D
MOP582\_6m, the practical impact is approximately $1^\circ$.

\subsection{Spectral shape}

Per-segment surface elevation spectra agree closely in shape.
Figure~\ref{fig:ruby2d_spec} compares both pipelines at a representative
median segment ($H_s = 0.77$~m) and at a storm peak
($H_s = 2.78$~m, bimodal sea state). For the spectrum overlay, the
17-minute multi-taper spectra were averaged over the three or four segments inside
the prior pipeline's 60-minute window so that both estimates cover the
same time period. Integrated $H_s$ from the spectra agrees to within
4\% in both cases. The visible difference is the chi-square noise of a
single-periodogram estimator versus the variance reduction from
averaging seven DPSS tapers; the underlying peak shapes and locations
agree.

\begin{figure}[h]
  \centering
  \includegraphics[width=0.95\textwidth]{Ruby2D/ruby2d_582_6m_spectra.png}
  \caption{Surface elevation spectra at two representative
    segments for Ruby2D MOP582\_6m. Top: median segment
    ($H_s = 0.77$~m, November 27, 2021). Bottom: storm peak
    ($H_s = 2.78$~m, December 14, 2021). Left: log scale, full SS
    band. Right: linear scale, swell band detail. The multi-taper
    estimates (blue) are smoother than the prior single-periodogram
    estimates (red); peak locations, peak heights, and total variance
    agree.}
  \label{fig:ruby2d_spec}
\end{figure}

\subsection{Implications}

Two conclusions follow from the head-to-head test:

\begin{itemize}
  \item The multi-taper pipeline reproduces an independent prior
    pipeline to within instrument noise on bulk wave parameters.
    Prior Ruby2D bulk wave statistics --- and any analyses that rely
    on them --- do not require revision.
  \item The improvements that multi-taper provides are exactly the
    expected ones: variance reduction in the spectral estimate and
    the elimination of edge cases (the prior pipeline's
    non-band-limited peak picker and inconsistent auto/cross
    windowing). Bulk parameters were already converged.
\end{itemize}

The Ruby2D test also surfaced one open question worth flagging: the
median Z-test (pressure-velocity consistency) ratio in the SS band
is 1.02 in the prior pipeline and 0.76 in the new pipeline, the same systematic
0.79 we observe across TBR23. The two pipelines use the same
definition, so the most likely explanations are that the prior
pipeline computes the diagnostic on uncorrected $S_{pp}$ rather than
the pressure-corrected $S_{\eta}$, or that we have a normalization
issue worth checking. This is on the followup list but does not
affect any of the bulk parameters reported above.

\section{Conclusions}

The full multi-taper ($NW = 4$, \texttt{nfft}~=~2048) should be the
default spectral method. The implementation is now a drop-in replacement
controlled by a single option
(\texttt{opts.spectralMethod~=~'mtm\_full'}). Backward compatibility
is preserved since $H_s$, $E_f$, and $S_{xx}$ are unchanged to within
0.1\%. $T_p$ improves from quantized 0.008~Hz steps to continuous
0.001~Hz resolution. Spectral shape diagnostics ($Q_p$, bandwidth,
spectral excess) are more physically meaningful.

The hybrid approach (DPSS within Welch sub-segments) does not offer
clear advantages at \texttt{nfft}~=~256 and introduces a spurious
$T_p$ shift and thus is not recommended.

Processing time increases from 0.7~min to 2.7~min per instrument
(TBR23 MOP580\_7m, 5285 segments). For the full pipeline (33
instruments), this adds $\sim$1 hour to a complete reprocessing run.

\begin{figure}[p]
  \centering
  \includegraphics[width=\textwidth]{spectral_method_comparison_spectra.png}
  \caption{Time-averaged surface elevation spectra for TBR23
    MOP580\_7m (8.3~m depth, 4795 segments). Left: full SS band.
    Right: swell peak detail. The full multi-taper (green) resolves a
    sharp swell peak that Welch (blue) smears across 3 frequency
    bins.}
  \label{fig:spectra}
\end{figure}

\begin{figure}[p]
  \centering
  \includegraphics[width=\textwidth]{spectral_method_comparison_bulk.png}
  \caption{Bulk parameter comparison across 4795 segments. Top row:
    scatter plots of MTM vs.\ Welch for $H_s$, $T_p$, directional
    spread, and $S_{xx}$. Bottom row: difference histograms for $H_s$
    and $T_p$. $H_s$ is identical across methods; $T_p$ shows
    discrete quantization bands in Welch (horizontal stripes in
    scatter) that the full multi-taper resolves continuously.}
  \label{fig:bulk}
\end{figure}

\begin{figure}[p]
  \centering
  \includegraphics[width=\textwidth]{spectral_method_comparison_single_seg.png}
  \caption{Single-segment spectral comparison at the peak energy
    segment ($H_s = 1.50$~m). The full multi-taper resolves a swell
    peak 2--3$\times$ taller than the Welch estimate. Total energy
    (area under the curve) is identical.}
  \label{fig:single}
\end{figure}

\begin{figure}[p]
  \centering
  \includegraphics[width=\textwidth]{spectral_method_comparison_ztest.png}
  \caption{Z-test distributions (pressure--velocity consistency) for
    the SS and IG bands. All three methods produce indistinguishable
    distributions, confirming the systematic departure from 1.0 is a
    data property, not a windowing artifact.}
  \label{fig:ztest}
\end{figure}

\end{document}
